\item Dos radiotransmisores $A$ y $B$ están separados por $10.0\text{ m}$. El transmisor $A$ opera a $60.0\text{ MHz}$, y el transmisor $B$ emite coherentemente en la misma frecuencia pero con una fase desfasada en $180^\circ$ respecto a $A$. ¿Qué distancia mínima debe recorrer un observador al caminar desde $A$ hacia $B$ sobre la línea que los une para encontrar el primer punto donde las dos señales estén exactamente en fase?